\subsubsection{Smart Order Routing}

\begin{definition} \hlt{Execution-Side SOR}\\
In execution, a smart order router chooses where each child order should go after the parent algorithm has chosen when and how much to trade. It is the venue-optimisation layer inside the execution stack.
\end{definition}

\begin{taxonomy} \hlt{Routing Logic Segmentation}
\begin{enumerate}[label=\roman*.]
\setlength{\itemsep}{0pt}
\item \hlt{observable liquidity}: quotes, depth, queue position, hidden-liquidity signals, and auction state;
\item \hlt{fill probability}: venue hit rate, cancellation speed, adverse selection, and reject history;
\item \hlt{net economics}: explicit fees, maker-taker rebates, clearing, and broker-specific schedules;
\item \hlt{latency}: feed delay, gateway delay, matching-engine delay, and stale-quote risk;
\item \hlt{information cost}: signalling from sweeping, probing, or repeatedly visiting the same venue.
\end{enumerate}
\end{taxonomy}

\begin{process} \hlt{Child-Order Routing Loop}
\begin{enumerate}[label=\roman*.]
\setlength{\itemsep}{0pt}
\item receive the child quantity, side, limit, and urgency from the execution algorithm;
\item score venues by expected net cost and completion probability;
\item split or sequence messages across venues;
\item update venue scores from fills, partial fills, rejects, and missed liquidity.
\end{enumerate}
\end{process}

\begin{caveat} \hlt{Displayed Best Price Is Not Enough}\\
A venue with the best quote can still be a poor destination if the queue is unstable, the quote is stale, fees are high, or fills are followed by adverse price moves.
\end{caveat}
